% Chapter 6: Conclusion
\chapter{Conclusion}
\label{ch:conclusion}

This chapter summarizes the research contributions of MutBench, analyzes its limitations, and presents future research directions.

\section{Summary of Research Contributions}
\label{sec:contribution_summary}

This dissertation presented MutBench, a systematic experimental framework for viral mutation hotspot detection.
The research contributions of MutBench are summarized in the following three points.

\textbf{Contribution 1. MutBench: an experimental framework for analyzing which factors influence viral mutation hotspot detection performance across 11 RNA viruses.}
MutBench introduced the first systematic experimental framework for viral mutation hotspot detection, designed to answer the question: \emph{which biological information source is most effective for which pathogen?}
It organizes 10 information types into 20 scoring formulas across 6 categories (frequency, MSA-derived, phylogenetic, structural, AI-based, and composite), defines a three-layer ground truth framework (Layer~A as a literature-curated, pathogen-specific union of convergent-evolution, immune-escape, and HVR-membership evidence types---with immune-escape sites dominating; Layer~B as purifying-selection negatives; Layer~C as DMS experimental fitness), introduces the hotspot-score composite metric, and presents a two-stage experimental design combining in-depth SARS-CoV-2 analysis (Stage~1) with large-scale comparison across 11 RNA viruses (Stage~2; see Chapter~\ref{ch:results} for detailed statistics).

\textbf{Contribution 2. Statistical demonstration that the optimal biological information source differs fundamentally across pathogens ($\omega^2_{\text{scoring} \times \text{pathogen}} = 0.296$), with 9 distinct optimal information types across 11 pathogens; LOPO 0/11 is reported as null-consistent corroboration of the interaction effect, not as independent evidence of non-generalizability.}
Statistical analysis of the large-scale benchmark results across 11 RNA viruses provided consistent evidence that the optimal biological information source is pathogen-dependent: the ANOVA scoring~$\times$~pathogen interaction and two null-consistent corroborations (LOPO 0/11 and Friedman) were computed on the Stage~2/3 8{,}580-evaluation dataset, and the vaccine-escape enrichment was computed on a separate 2{,}340-evaluation 3-pathogen subset (HIV-1, H3N2, SARS-CoV-2) with curated antibody-escape annotations.
First, the scoring $\times$ pathogen interaction ($\omega^2 = 0.296$ at 20-type granularity; 95\% cluster-bootstrap CI [0.195, 0.333]; $\omega^2 = 0.103$ at 6-category granularity as a collinearity-robust lower bound) was the largest identified source of variance among modeled terms in the factorial ANOVA, with within-cell residual ($\omega^2 \approx 0.39$) representing unmodeled variation including the three-way interaction (Chapter~\ref{ch:results} Table~\ref{tab:stage2_anova}).
Second, all 11 pathogens exhibited unique best scoring--detection combinations (11/11 unique best combos), with 9 distinct optimal information types, yielding a mean MCC gap of 0.265 between oracle and generalized performance.
Third, the Friedman test ($p = 0.990$) showed no reliably superior combination across all pathogens, consistent with pathogen-dependent optimality.
Fourth, Leave-One-Pathogen-Out (LOPO) cross-validation yielded a 0/11 agreement result. Because this outcome is null-consistent with random-combination permutation, we treat it as a corroboration of the ANOVA interaction rather than as independent evidence against generalization (the oracle-vs-generalized MCC gap of 0.265 is the quantitative anchor).

\textbf{Contribution 3 (two-part).} \textit{(3a) Multi-source integration infrastructure}: a compact 4-feature core (homoplasy, pLDDT, entropy, frequency) retains $\approx$98\% (97.6\%) of the full 10-feature EqualWeight ensemble mean MCC on the 12-pathogen Stage~3 panel under uniform-weight LOPO top-10\% evaluation (0.081 vs.\ 0.083; numerical equivalence, no paired-test significance), and EqualWeight integration achieves a statistically significant 4.012-fold H3N2 enrichment ($p = 0.0023$) where single-feature top-5 approaches (FreqThresh top-5, BestSingle top-5) fail to reach significance in the same integration-block comparison---this scope-restricted result does not contradict the existence of single-scoring combinations that survive Bonferroni correction within the full 780-combination Stage~2 sweep on H3N2 (Section~\ref{sec:cross_validation_escape}). \textit{(3b) Externally-valid vaccine-escape enrichment, anchored on HIV-1}: the Layer-A-disjoint subset of HIV-1 escape positions (37 of 45) yields a direct novel-only enrichment of $7.16\text{--}8.24\times$ (worst-case nominal Fisher $p = 5.0 \times 10^{-12}$, Bonferroni-adjusted across 780 combinations $p_{\text{adj}} \approx 3.9 \times 10^{-9}$); H3N2's 9.36$\times$ is a self-consistency check (69\% Layer~A overlap); SARS-CoV-2 is exploratory after multiple-comparison correction.
A systematic analysis of 10 biological information types (Section~\ref{sec:information_analysis}) reveals that the most informative feature differs across pathogens---homoplasy for Norovirus (AUC\,=\,0.944), pLDDT for H3N2 (0.787), ESM-2 LLR for SARS-CoV-2 (0.586)---and that feature ablation identifies a 4-feature core (homoplasy, pLDDT, entropy, frequency) capturing $\approx$98\% (97.6\%) of integrated performance.
Multi-source integration reduces the experimental search space, achieving H3N2 vaccine escape enrichment of 4.012$\times$ ($p = 0.0023$) where single-information approaches fail to reach significance.
Vaccine-escape enrichment analysis across 3 of 11 pathogens (H3N2, SARS-CoV-2, HIV-1; 2{,}340 evaluations; selected for availability of curated antibody-escape annotations) was audited for circularity against the Layer~A ground truth. \textbf{HIV-1 is the cleanest external validation}: its escape set shares only 8 of 45 positions (18\%) with HIV-1 Layer~A, so the 7.19$\times$ enrichment (Bonferroni-adjusted $p = 2.5 \times 10^{-16}$) primarily reflects detection of positions outside the benchmark's training signal. \textbf{H3N2 should be read as a self-consistency check, not external validation}: 20 of 29 H3N2 escape positions (69\%) are in Layer~A, and on the 9 Layer~A-disjoint (novel) escape positions the top combination's enrichment drops to 7.19$\times$ with Fisher $p = 0.132$ (not significant under any correction), so the headline H3N2 9.36$\times$ predominantly re-measures Layer~A performance under a different label rather than discovering new escape positions. \textbf{SARS-CoV-2 is exploratory only}: the 7.63$\times$ nominal enrichment does not survive Bonferroni correction ($p_{\text{adj}} \to 1$), and 60\% of its escape set overlaps Layer~A. The practical-impact claim of Contribution~3 therefore rests on HIV-1 as the primary external validation, with H3N2 and SARS-CoV-2 contributing self-consistency and exploratory evidence respectively.
This analysis establishes which biological information is essential for each pathogen and provides a ranked candidate list for experimental validation.


\section{Limitations}
\label{sec:limitations}

Table~\ref{tab:limitations_summary} summarizes the key limitations of MutBench and their current mitigation status.

\begin{table}[htbp]
\centering
\caption{Summary of MutBench limitations and mitigation status}
\label{tab:limitations_summary}
\small
\adjustbox{max width=\textwidth}{
\begin{tabular}{llll}
\toprule
\textbf{Limitation} & \textbf{Severity} & \textbf{Current mitigation} & \textbf{Future work} \\
\midrule
Ground truth limited to literature & High & 3-layer GT with DMS validation & Expand to more pathogens \\
DMS data available for 6/11 pathogens & Medium & Layers A+B for remaining 5 & New DMS experiments (e.g., H5~\cite{dadonaite2024h5}) \\
Phylogenetic non-independence & High & PoC simulation + TreeTime validation & Full pipeline integration \\
Synthetic--real performance gap & Medium & Dual benchmark design & Adaptive reference sequences \\
Position-level exact matching & Medium & Enrichment ratio as complement & Region-overlap metrics \\
RNA viruses only (11 pathogens) & Low & Covers major RNA families & DNA virus extension \\
GISAID sampling bias & Medium & Weekly temporal stratification & Geographic stratification \\
Single protein per pathogen & Low & Surface protein is primary target & Multi-protein evaluation \\
GT completeness (false negatives) & Medium & Literature-curated; DMS cross-validation & Expand Layer A with new studies \\
MCC cross-pathogen comparability & Low & ANOVA uses relative rankings & Prevalence-adjusted metrics \\
Layer~A $\cap$ escape overlap (H3N2 69\%, SARS 60\%) & High & HIV-1 (18\%) as primary external anchor; novel-only audit reported & Additional pathogens with Layer~A-disjoint escape sets \\
$\rho$(Layer~A, frequency) pathogen-specific (mean $|\rho|{=}0.12$; HCV 0.33, Inf.B 0.26, Norovirus 0.23) & Medium & Homoplasy/DMS as frequency-independent signals; Ch5 circularity audit & Uniform Layer~A curation protocol across pathogens \\
Winner's curse in top-1 enrichment & Medium & Bonferroni-survivor mean reported ($\approx 4.7\times$ vs top 9.36$\times$) & Bootstrap-shrunken point estimates \\
Technical covariates (length, pos.\ rate) $\to$ MCC & Medium & Ch5 sensitivity: $\rho = -0.73$/$+0.64$ disclosed; ANCOVA-adjusted scoring~$\times$~pathogen $\omega^2$ rises 0.234~$\to$~0.264 (Section~\ref{sec:cross_pathogen_implications}) & Matched-covariate resampling on disjoint pathogen panel \\
LOPO 0/11 expected under random null ($P \approx 0.96$) & Medium & Permutation-null analysis in Chapter~\ref{ch:results} (Section~\ref{subsec:9pathogen_lopo}); gap 0.265 as primary evidence & Larger pathogen panel (20--30+) for power \\
Vaccine-escape available for 3/11 pathogens & Medium & Ch4 Table \ref{tab:escape_availability} availability audit & Curation effort on Dengue/RSV/Rabies escape sets \\
Head-to-head comparison with concurrent benchmarks (ViroGym~\cite{virogym2026}, EVEREST~v3~\cite{gurev2026everest_v3}, PLANT~\cite{plant2025antigenic}) not performed & Medium & Task-orthogonality is structural rather than a deferral: ViroGym/EVEREST optimise per-variant fitness Spearman, MutBench optimises per-region MCC; the metrics do not share a common evaluation universe even on overlapping pathogens. The pathogen-overlap subset (HIV-1, SARS-CoV-2, H3N2 are in both panels) is reported in this work as the vaccine-escape cross-validation result, not as a model-vs-model leaderboard. & A formal task-aligned subset leaderboard (e.g., binarising fitness predictions at top-$k$ for region-level evaluation) is queued for the defense-week sprint. \\
\bottomrule
\end{tabular}
}
\end{table}

\textbf{Ground Truth Limitations.}

The full three-layer ground truth (Layers A, B, and C) is available only for the six pathogens with DMS data (SARS-CoV-2, H3N2, HIV-1, RSV, Rabies, EV-A71); the remaining five pathogens are evaluated using only Layers A and B.
Layer~C (DMS functional fitness) is based on published DMS datasets~\cite{dadonaite2023dms}, and for pathogens without available DMS data, functional fitness validation cannot be performed.
Furthermore, in interpreting DMS data, the choice of threshold for defining functionally important positions (top 20\% of mean absolute fitness effect in this study) directly affects the composition of the ground truth, introducing sensitivity to that threshold.
DMS measures in vitro molecular fitness using pseudovirus systems and thus may not fully capture in vivo fitness.

\textbf{Methodological Limitations.}

The phylogenetic non-independence of frequency-based scores (Section~\ref{sec:phylo_limits}) and GISAID sampling biases are discussed in Chapter~\ref{ch:discussion}.
Additionally, the current implementation of MutBench targets only point mutations, excluding insertions and deletions (indels), and validation on DNA viruses has not yet been performed.
A performance gap exists between synthetic and real data due to the sparse distribution of real H-scores, though biological validity is preserved with enrichment ratios above 6.9-fold (see Chapter~\ref{ch:discussion} for detailed analysis).

\textbf{Limitations of Position-Level Evaluation.}

Because MutBench computes precision and recall at the exact nucleotide position level, a detection that is off by even a single position from the ground truth is classified as a false positive.
Window-based or region overlap metrics could more leniently reflect the practical utility of detection.

\section{Future Research Directions}
\label{sec:future_work}

Table~\ref{tab:future_roadmap} presents the prioritized future research roadmap, organized by expected timeline and impact.

\begin{table}[htbp]
\centering
\caption{Prioritized future research roadmap}
\label{tab:future_roadmap}
\small
\adjustbox{max width=\textwidth}{
\begin{tabular}{llll}
\toprule
\textbf{Priority} & \textbf{Research direction} & \textbf{Prerequisite} & \textbf{Expected impact} \\
\midrule
1 (Highest) & Phylogenetic correction integration & TreeTime pipeline & Resolves founder-effect bias \\
2 & Adaptive method selection & Corrected benchmark + 20+ pathogens & Pathogen-specific optimization \\
3 & DNA virus extension & New ground truth curation & Broader applicability \\
4 & Real-time surveillance integration & Incremental H-score updates & Public health utility \\
5 & Independent multi-tool reproduction of the 8{,}580-evaluation benchmark on the same pathogen panel (Bayesian partial-pooling, ART ANOVA, ANCOVA, $\alpha = 0.01$ sensitivity, and HIV-1 novel-only Fisher $p$ are reported in Section~\ref{sec:cross_pathogen_implications} of Ch5; only independent reproduction remains open) & Existing 8{,}580-evaluation result archive & Closes the single-team / single-codebase reproducibility gap \\
6 & Supplementary-table extensions of the existing provenance manifest (the row-to-cell manifest itself is already bundled; remaining items: DMS preprocessing convention table; region-overlap LOPO mini-table; Random-Forest CV AUC summary; Bonferroni-survivor median-of-survivors column; Tajima's D / CUSUM / segmentation citation reconciliation) & Existing per-pathogen CSV archive + bib audit & Closes the cycle-12 P2 polish set in one supplementary-table sweep \\
\bottomrule
\end{tabular}
}
\end{table}

\textbf{Adaptive Method Selection.}

The oracle ceiling analysis (Section~\ref{subsec:adaptive_weighting_limits}) demonstrated that per-pathogen feature weighting can double performance (MCC = 0.160 vs.\ EqualWeight 0.083), but with only 11 reference pathogens the knowledge base is insufficient for reliable adaptive selection.
Scaling to 20+ pathogens, introducing top-$k$ ensemble strategies, and transitioning to meta-learning regression models~\cite{rice1976algorithm,kerschke2019algorithm} are the concrete next steps toward pathogen-adaptive method selection.

\textbf{Integration of Phylogenetic Correction.}

The proof-of-concept results reported in Section~\ref{sec:phylo_simulation} establish that homoplasy-based scoring captures qualitatively different evolutionary signals from frequency-based H-scores: $\rho = -0.876$ in coalescent simulation, AUC~$= 1.000$ for TreeTime ML ancestral reconstruction ($N = 2{,}000$), and $\rho = -0.107$ on real GISAID sequences ($N = 500$).
A concrete integration pathway would replace H-score with homoplasy count---the number of independent mutational origins inferred by ancestral reconstruction---as the input to scoring formulas, requiring TreeTime~\cite{sagulenko2018treetime} ancestral reconstruction as a preprocessing step for each pathogen.
The primary computational bottleneck is TreeTime's $O(N \log N)$ scaling with sequence count, but recent ultrafast alternatives such as UShER~\cite{turakhia2020ultrafast} enable processing of millions of sequences, making this integration feasible at GISAID scale.
Phylogenetic correction would primarily affect pathogens with strong founder effects (e.g., SARS-CoV-2 D614G, which ranks 2nd by H-score but 1,134th by homoplasy), potentially changing the best-performing scoring--detection combinations for those pathogens and improving the overall oracle MCC.
Temporal stratification, which normalizes sequence proportions across time periods, would complement phylogenetic correction by mitigating sequencing bias independently of tree topology.

\textbf{Additional Research Directions.}

Additional research directions include:
(1)~expansion to DNA viruses (HPV, HBV) whose lower mutation rates and distinct mutation mechanisms (e.g., APOBEC3-mediated deamination) require verification of current scoring and detection applicability;
(2)~indel handling and improved PLM-derived features (the current 20-scoring panel already includes ESM-2 LLR~\cite{lin2023esm2,meier2021language} and an ESM-2 semantic-change channel; ``improved'' here means alternative or virus-fine-tuned PLMs such as ProtTrans~\cite{elnaggar2022prottrans} and the multimodal ESM-3 generation, plus principled gap-handling strategies for indels that are presently excluded from the H-score framework);
and (3)~real-time surveillance pipeline development, integrating MutBench-validated detection methods with incremental H-score updates for streaming sequencing data and automatic alerts upon novel variant emergence.

\section{Concluding Remarks}
\label{sec:final_conclusion}

The COVID-19 pandemic revealed that viral mutation hotspot detection is not merely an academic exercise but a public health imperative.
This work demonstrates that the most impactful factor in hotspot detection is selecting the appropriate biological information source for each pathogen. The recommended scoring categories by virus family (Table~\ref{tab:scoring_recommendation}) provide practical guidance for researchers analyzing new viral sequences.

This dissertation presented MutBench, an experimental framework that enabled this finding by systematically comparing 10 biological information types across 11 RNA viruses (8{,}580 total evaluations).
The central result---that the scoring $\times$ pathogen interaction ($\omega^2 = 0.296$; 11-pathogen cluster-bootstrap 95\% interval [0.195, 0.333]) is the largest identified modeled variance component, alongside a within-cell residual $\omega^2 \approx 0.39$ that reflects further pathogen-level heterogeneity---empirically supports the algorithm selection principle~\cite{kerschke2019algorithm} and provides evidence that the optimal biological information source is pathogen-dependent.
Nine distinct optimal information types were identified across 11 pathogens, from FUBAR phylogenetic selection scores for H3N2 to protein language models for SARS-CoV-2, indicating that no single information source suffices across the viral landscape.
The concordance between Layer~A benchmark rankings and vaccine escape enrichment is strongest for HIV-1 (7.19$\times$, 82\% of the escape set disjoint from Layer~A; Bonferroni-adjusted $p = 2.5 \times 10^{-16}$), which provides the cleanest external validation; H3N2's headline 9.36$\times$ is largely a self-consistency check given 69\% Layer~A overlap, and SARS-CoV-2's 7.63$\times$ is exploratory (Bonferroni-adjusted $p \to 1$). The benchmark's practical utility for identifying immune evasion positions is therefore established with high confidence only on the Layer~A-disjoint subset exemplified by HIV-1.
Multi-source integration further reduces the experimental search space, with a compact 4-feature core capturing $\approx$98\% (97.6\%) of full-ensemble performance.
These conclusions are bounded by limitations enumerated in Section~\ref{sec:limitations}: the eleven-pathogen scope precludes frequentist mixed-effects modeling (a Bayesian partial-pooling alternative is reported in Section~\ref{sec:cross_pathogen_implications}); the covariate-adjusted (length, positive-rate) variance decomposition is performed in the same section but on the same dataset; the LOPO 0/11 result is null-consistent under random-combination permutation; the Layer~A construction is partly literature-curated and dominated by immune-escape annotations; the multi-source enrichment claim is statistically anchored by HIV-1; the scoring panel is single-position and does not model epistasis or coupled substitutions; and head-to-head comparison with concurrent benchmarks (ViroGym~\cite{virogym2026}, EVEREST~v3~\cite{gurev2026everest_v3}) remains future work.
As new viral threats emerge, the information-type comparison methodology established here enables rapid identification of optimal information sources tailored to each pathogen's characteristics.
\textbf{Take-home message}: within this single-position scoring regime and the modeled-effect terms of the factorial ANOVA, choosing the right biological information source for the pathogen at hand explains more performance variation than choosing among detection algorithms---and quantifying that asymmetry is what MutBench contributes.

\paragraph{Dual-use considerations.} MutBench predicts positions that are likely to mutate under selection, including immune-escape-relevant sites on viral surface glycoproteins. The same outputs that inform vaccine-strain selection, surveillance triage, and antibody design can in principle inform engineering of escape variants, so the benchmark outputs are inherently dual-use. Three mitigations were therefore built into the release plan rather than being deferred to downstream users: (i)~raw GISAID sequences are never redistributed---only computed H-scores, detection results, and summary statistics are shared, preserving the provider-side access controls of the underlying data agreement; (ii)~the benchmark publishes \emph{region-level} hotspot rankings, not variant-level engineering recipes, so it raises the signal-to-noise for surveillance without providing a direct uplift to escape-variant design; and (iii)~the public repository release (30~days post-defence) coincides with disclosure to the relevant biosafety review body, and any extension of the benchmark to pathogens on the U.S.\ HHS P3CO/Select Agent lists will go through a dedicated institutional biosafety review prior to publication. These steps do not eliminate dual-use risk, but they keep the benchmark aligned with the \emph{responsible stewardship} expectations now standard for computational pathogen-evolution work.

\paragraph{Code and data availability.}\label{sec:code_availability}
The source code of the MutBench framework, benchmark scripts, ground truth definitions for 11 pathogens, and configuration files necessary for reproducing all experiments are archived at the GitHub repository (\url{https://github.com/hwijunkwon/mutbench}; the submission-time state is captured by tag \texttt{dissertation-v1}, SHA fingerprint to be inserted at deposit).%
\footnote{Repository access during the review period: the tagged snapshot is available to examiners via the supplementary materials bundle; a public release mirrored on Zenodo (DOI to be assigned at deposit) is scheduled within 30~days of the defense.}
 All detection method implementations, ground truth definitions, and evaluation scripts are organized as a reproducible pipeline with configuration files specifying exact parameters for each experiment.
The benchmark can be reproduced with three commands: (1)~feature computation (\texttt{python compute\_features.py --pathogen H3N2}), (2)~detection evaluation (\texttt{python run\_stage3\_full\_detectors.py}), and (3)~statistical analysis (\texttt{python analyze\_stage3\_stats.py}). The full 8{,}580-evaluation benchmark (20 scoring types $\times$ 39 detectors $\times$ 11 pathogens) completes in 25 seconds on a single workstation (measured: 24.96s on an Intel i7 with RTX 4090, without GPU dependency for the benchmark itself). Feature computation for a new pathogen requires approximately 15--20 minutes including MSA construction, frequency/entropy calculation, homoplasy inference, and ESM-2 scoring; structural features (pLDDT, SASA) are optional and require an AlphaFold or experimental structure when available.
The repository becomes publicly accessible within 30~days of the dissertation defense, with the submission-time tag \texttt{dissertation-v1} and the matching Zenodo archive serving as the immutable reference for all numbers reported in this thesis.
The benchmark was developed using Python~3.10 with NumPy~1.24, SciPy~1.11, Pandas~2.0, Matplotlib~3.7, and Seaborn~0.12.
MSA construction used MAFFT~v7.505 with the \texttt{--auto} strategy.
Phylogenetic trees were inferred using IQ-TREE~2 with ModelFinder; selection pressure analysis used HyPhy~v2.5 (FUBAR, MEME); protein language model scores used the \texttt{esm2\_t33\_650M\_UR50D} checkpoint; and the \texttt{stability} scoring channel uses a BLOSUM62 substitution proxy (see Section~\ref{sec:scoring_detection}), chosen in place of physics-based $\Delta\Delta G$ predictors such as FoldX that are computationally infeasible at the 11-pathogen $\times$ 39-detector scale.
All stochastic processes (bootstrap resampling, subsampling) used a fixed random seed (42) for reproducibility.
Pathogen sequence data were collected from NCBI GenBank, and SARS-CoV-2 DMS data were obtained from the original data provided by Dadonaite et al.~\cite{dadonaite2023dms, dadonaite2024nature}.
GISAID sequence data were accessed under the GISAID Data Access Agreement. The benchmark framework does not redistribute raw GISAID sequences; instead, computed H-scores and detection results are provided.

\paragraph{Numerical artifacts and CSV-to-table provenance.}
Core Stage~3 and vaccine-escape numeric claims in Chapter~\ref{ch:results} resolve to specific CSVs in the \texttt{results/} directory of the archived repository (a small set of remaining provenance gaps---primarily Stage~1 real-data and adaptive-weighting CSVs---is enumerated in the bundled provenance manifest). The primary artifacts are \texttt{stage3\_full\_results.csv} (8{,}580-row Stage~2 evaluation output; source of Tables~\ref{tab:stage2_best} and \ref{tab:stage2_anova}), \texttt{layer\_c\_evaluation.csv} (Layer~C DMS cross-check), \texttt{feature\_analysis/per\_feature\_auc\_11pathogens.csv}, \texttt{feature\_analysis/rf\_cv\_auc\_v2.csv} (Table~\ref{tab:rf_cv_auc}), and \texttt{vaccine\_escape\_stage3.csv} (single-scoring block of Table~\ref{tab:vaccine_escape_stage3}), and \texttt{escape\_validation\_adapt.csv} (EqualWeight integration block of the same table, \texttt{threshold=top5}). These files together with a CSV-row-to-table-cell provenance manifest are bundled in the dissertation examiners' submission archive and mirrored alongside the Zenodo deposit, so each reported number in the thesis can be re-derived by reading the listed CSV row rather than by re-running the pipeline.
